%═══════════════════════════════════════════════════════════════════════════════
\chapter{Sobreajuste, regularización y generalización}
\label{ch:cap2}
%═══════════════════════════════════════════════════════════════════════════════

El Capítulo~\ref{ch:cap1} cerró con una promesa: existe una familia paramétrica
no lineal capaz de aproximar cualquier función continua con precisión arbitraria.
Ahora surge la pregunta inmediata: si podemos aproximar cualquier función, ¿por
qué no interpolar exactamente los datos de que disponemos? La respuesta define uno
de los conceptos centrales del aprendizaje automático, y es, en el fondo, la misma
advertencia que el análisis numérico lleva décadas haciendo sobre la interpolación
de alto grado.

Este capítulo construye el marco conceptual que responde esa pregunta con precisión.
El recorrido va de la observación numérica al análisis espectral: primero se muestra
el problema con un ejemplo concreto, luego se formaliza la distinción entre error de
entrenamiento y error de generalización, después se analiza el dilema sesgo-varianza,
y finalmente se desarrolla la regularización de Tikhonov, su interpretación como
filtro espectral mediante la SVD, y los criterios para seleccionar el parámetro de
regularización.

%───────────────────────────────────────────────────────────────────────────────
\section{El peligro de la interpolación perfecta}
\label{sec:runge}
%───────────────────────────────────────────────────────────────────────────────

\dificultad{2}{2}

El estudiante de análisis numérico conoce el fenómeno de Runge: un polinomio de
grado alto que interpola exactamente una función suave en nodos equiespaciados puede
exhibir oscilaciones violentas entre los nodos, especialmente cerca de los extremos
del intervalo. Este fenómeno no es una curiosidad numérica; es la primera manifestación
de un problema estructural que reaparece en todas las formas de aprendizaje: el modelo
que memoriza los datos no necesariamente describe bien el fenómeno subyacente.

Para verlo en el lenguaje del Capítulo~\ref{ch:cap1}, considérese la función de
Runge:
\[
  g(x) = \frac{1}{1 + 25x^2}, \qquad x \in [-1, 1],
\]
y un conjunto de $N$ nodos equiespaciados $x_i = -1 + 2(i-1)/(N-1)$ con valores
$y_i = g(x_i)$. La familia polinomial de grado $d = N-1$ interpola exactamente:
$\mathcal{L}_{\text{train}}(\bm{\theta}^*) = 0$. Sin embargo, al evaluar el polinomio
en puntos intermedios, el error puede ser enorme. Para $N = 11$, el error máximo
del polinomio interpolante de grado 10 supera el 100\% del rango de la función.

El mensaje es preciso:

\begin{observacion}{Pérdida empírica cero no implica buen modelo}
Minimizar la pérdida empírica $\mathcal{L}_{\text{train}}(\bm{\theta})$ hasta cero
equivale a exigir que el modelo pase exactamente por todos los puntos de datos.
Cuando el número de parámetros es igual al número de datos, esto siempre es posible.
Pero el modelo resultante describe el ruido de las observaciones, no la función
subyacente. Este fenómeno se denomina \textbf{sobreajuste} (\textit{overfitting}).
\end{observacion}

\begin{fisicaconexion}{Memorizar datos vs.\ describir la naturaleza}
El fenómeno de Runge es la versión matemática de un problema experimental que todo
físico ha encontrado en el laboratorio: cuando se ajusta un polinomio de grado alto
a datos con ruido, el ajuste puede pasar perfectamente por cada punto medido pero
oscilar de forma incontrolada entre ellos. Un ejemplo inmediato es la deconvolución
espectral: si se intenta reconstruir el perfil de una línea espectral ajustando un
polinomio interpolante a los canales del espectrómetro, el resultado suele oscilar
entre canales y producir ``intensidades negativas'' entre los máximos, que son
físicamente imposibles. El problema no está en los datos ni en el método: está en
que el modelo tiene demasiada libertad para absorber el ruido de cada medición
individual en lugar de la señal subyacente. La regularización es el mecanismo
que impone la suavidad física de la señal como restricción matemática.
\end{fisicaconexion}

El sobreajuste tiene una causa matemática identificable: la familia paramétrica
tiene demasiados grados de libertad en relación con la información contenida en los
datos. Formalizar esta idea requiere distinguir entre dos tipos de error.

%───────────────────────────────────────────────────────────────────────────────
\section{Error de entrenamiento y error de generalización}
\label{sec:train_test}
%───────────────────────────────────────────────────────────────────────────────

\dificultad{1}{1}

\begin{definicion}{Partición de datos: entrenamiento y evaluación}
\label{def:particion}
Sea $\mathcal{D} = \{(\bm{x}_i, y_i)\}_{i=1}^N$ un conjunto de datos. Una
\textbf{partición entrenamiento-evaluación} con fracción $\alpha \in (0,1)$ consiste
en dividir $\mathcal{D}$ en dos subconjuntos disjuntos:
\[
  \mathcal{D}_{\text{train}} \;\cup\; \mathcal{D}_{\text{test}} = \mathcal{D},
  \qquad
  \mathcal{D}_{\text{train}} \;\cap\; \mathcal{D}_{\text{test}} = \emptyset,
\]
con $|\mathcal{D}_{\text{train}}| = \lfloor \alpha N \rfloor$ y
$|\mathcal{D}_{\text{test}}| = N - \lfloor \alpha N \rfloor$. El subconjunto
$\mathcal{D}_{\text{train}}$ se usa para determinar $\bm{\theta}^*$; el subconjunto
$\mathcal{D}_{\text{test}}$ no participa en ese proceso y se reserva para evaluar
el modelo resultante.
\end{definicion}

\begin{definicion}{Error de entrenamiento y error de generalización}
\label{def:errores}
Sea $\bm{\theta}^* = \argmin_{\bm{\theta}} \mathcal{L}_{\text{train}}(\bm{\theta})$
el minimizador de la pérdida sobre $\mathcal{D}_{\text{train}}$. Se definen:
\begin{align}
  E_{\text{train}} &= \mathcal{L}_{\text{train}}(\bm{\theta}^*)
  = \frac{1}{|\mathcal{D}_{\text{train}}|}
    \sum_{(\bm{x}_i,y_i)\in\mathcal{D}_{\text{train}}}
    \ell\bigl(f(\bm{x}_i;\bm{\theta}^*),\, y_i\bigr),
  \label{eq:etrain}\\[6pt]
  E_{\text{test}} &= \mathcal{L}_{\text{test}}(\bm{\theta}^*)
  = \frac{1}{|\mathcal{D}_{\text{test}}|}
    \sum_{(\bm{x}_i,y_i)\in\mathcal{D}_{\text{test}}}
    \ell\bigl(f(\bm{x}_i;\bm{\theta}^*),\, y_i\bigr).
  \label{eq:etest}
\end{align}
La brecha $E_{\text{test}} - E_{\text{train}} \geq 0$ se denomina
\textbf{brecha de generalización}. Cuando esta brecha es grande, el modelo
ha sobreajustado.
\end{definicion}

\begin{observacion}{El conjunto de test como árbitro imparcial}
La condición de que $\mathcal{D}_{\text{test}}$ no participe en la determinación
de $\bm{\theta}^*$ es fundamental. Si se usaran datos de test para ajustar el modelo
o seleccionar hiperparámetros, $E_{\text{test}}$ dejaría de ser una estimación
imparcial del error de generalización. Este principio se viola frecuentemente en
la práctica y es fuente de estimaciones de rendimiento demasiado optimistas.
\end{observacion}

\subsection{La curva de capacidad}
\label{subsec:curva_capacidad}

\dificultad{2}{2}

La relación entre la complejidad del modelo y los errores de entrenamiento y test
tiene una estructura sistemática. Definimos la \textbf{capacidad} de una familia
paramétrica como el número de parámetros libres $p$. Al variar $p$ y registrar
$E_{\text{train}}$ y $E_{\text{test}}$, se obtiene la curva característica de la
Figura~\ref{fig:bias_variance_curve}.

\begin{figure}[H]
\centering
\begin{tikzpicture}[scale=1.0]
  % Axes
  \draw[->] (0,0) -- (8.5,0) node[right] {Capacidad $p$};
  \draw[->] (0,0) -- (0,5.5) node[above] {Error};

  % Training error curve (decreasing)
  \draw[thick, defcolor, domain=0.3:8.2, samples=120]
    plot (\x, {3.8*exp(-0.55*\x) + 0.15});

  % Test error curve (U-shape)
  \draw[thick, ejemcolor, domain=0.3:8.2, samples=120]
    plot (\x, {3.8*exp(-0.55*\x) + 0.15 + 1.1*exp(-1.6*(\x-4.5)^2)*(\x-1.5) + 0.05*\x});

  % Labels
  \node[defcolor, right] at (8.2, 0.55) {\small $E_{\text{train}}$};
  \node[ejemcolor, right] at (8.2, 1.25) {\small $E_{\text{test}}$};

  % Optimal zone
  \draw[dashed, gray] (4.5,0) -- (4.5,4.2);
  \node[below, gray] at (4.5,0) {\small óptimo};

  % Regions
  \node[defcolor!60, align=center] at (1.8, 4.5) {\small alto sesgo\\[-2pt]\small baja varianza};
  \node[ejemcolor!80, align=center] at (7.0, 4.5) {\small bajo sesgo\\[-2pt]\small alta varianza};
\end{tikzpicture}
\caption{Curva de capacidad. El error de entrenamiento decrece monótonamente al aumentar la capacidad $p$. El error de test tiene un mínimo en la zona de capacidad óptima; más allá de ese punto el modelo sobreajusta.}
\label{fig:bias_variance_curve}
\end{figure}

Cuatro hechos caracterizan esta curva, y todos tienen demostración:

\begin{enumerate}
  \item $E_{\text{train}}$ es no creciente en $p$: al añadir parámetros, el conjunto
        de modelos factibles se agranda y el mínimo de $\mathcal{L}_{\text{train}}$
        solo puede disminuir o mantenerse.
  \item Para $p = N$ (interpolación exacta), $E_{\text{train}} = 0$ con pérdida MSE
        y nodos distintos.
  \item $E_{\text{test}}$ tiene un mínimo para algún $p^* < N$ si los datos contienen
        ruido.
  \item La brecha $E_{\text{test}} - E_{\text{train}}$ es no decreciente en $p$.
\end{enumerate}

Los hechos 1 y 2 son consecuencias directas de la Proposición~\ref{prop:interpolacion}
del Capítulo~\ref{ch:cap1}. Los hechos 3 y 4 requieren el análisis sesgo-varianza
de la sección siguiente.

%───────────────────────────────────────────────────────────────────────────────
\section{La descomposición sesgo-varianza}
\label{sec:biasvar}
%───────────────────────────────────────────────────────────────────────────────

\dificultad{3}{3}

Para analizar la brecha de generalización con precisión necesitamos un modelo
del ruido. Se asume el modelo aditivo estándar:

\begin{definicion}{Modelo de datos con ruido aditivo}
\label{def:ruido}
Los datos se generan mediante:
\[
  y_i = g(\bm{x}_i) + \varepsilon_i,
  \qquad i = 1, \ldots, N,
\]
donde $g: \mathbb{R}^n \to \mathbb{R}$ es la \textbf{función objetivo} desconocida,
y $\varepsilon_i$ son perturbaciones aleatorias con:
\[
  \mathbb{E}[\varepsilon_i] = 0,
  \qquad
  \mathbb{E}[\varepsilon_i^2] = \sigma_\varepsilon^2,
  \qquad
  \mathbb{E}[\varepsilon_i \varepsilon_j] = 0 \text{ para } i \neq j.
\]
Es decir, el ruido tiene media cero, varianza $\sigma_\varepsilon^2$ y es
no correlacionado entre observaciones.
\end{definicion}

El operador $\mathbb{E}[\cdot]$ denota la \textbf{esperanza} sobre las
realizaciones del ruido $\bm{\varepsilon} = (\varepsilon_1,\ldots,\varepsilon_N)\T$:
para una función $h(\bm{\varepsilon})$, $\mathbb{E}[h(\bm{\varepsilon})] =
\int h(\bm{\varepsilon})\, p(\bm{\varepsilon})\, d\bm{\varepsilon}$, donde
$p(\bm{\varepsilon})$ es la densidad de probabilidad del ruido. En el modelo
anterior no se asume ninguna distribución específica; solo se usan las propiedades
de linealidad de $\mathbb{E}[\cdot]$ y las condiciones sobre los momentos del ruido.

\begin{observacion}{Por qué asumir ruido aditivo}
El modelo $y_i = g(\bm{x}_i) + \varepsilon_i$ es el más simple que captura la
realidad: las observaciones físicas siempre tienen alguna fuente de incertidumbre
(ruido instrumental, variabilidad natural, truncamiento numérico). La hipótesis de
media cero y varianza finita es suficiente para el análisis que sigue. La hipótesis
de no correlación es más fuerte, pero es la que corresponde a medidas independientes
en el sentido experimental usual.
\end{observacion}

\begin{fisicaconexion}{Sesgo y varianza como error sistemático y estadístico}
La distinción entre sesgo y varianza es exactamente la distinción entre
\textbf{error sistemático} y \textbf{error estadístico} que todo físico aprende
en su primer curso experimental. El error estadístico se reduce promediando más
mediciones: al aumentar $N$, la varianza del estimador disminuye como $1/N$.
El error sistemático, en cambio, no desaparece con más datos: si el instrumento
tiene un offset de calibración, la media de mil mediciones seguirá estando
desplazada del valor verdadero. En el contexto del ajuste de modelos, el sesgo
es el ``error de calibración'' del modelo: si la familia de funciones no puede
representar la función real $g$ (por ejemplo, se usa una recta para ajustar
una parábola), el estimador convergirá a un valor sistemáticamente equivocado
sin importar cuántos datos se tengan. La regularización introduce
deliberadamente un sesgo controlado para reducir el error estadístico (varianza):
es el mismo compromiso que un físico acepta cuando aplica un filtro suavizante
a sus datos.
\end{fisicaconexion}

\begin{definicion}{Sesgo y varianza de un estimador}
\label{def:sesgvar}
Sea $\hat{\bm{\theta}}(\mathcal{D}_{\text{train}})$ el estimador obtenido como
minimizador de $\mathcal{L}_{\text{train}}$, visto como función del conjunto de
datos de entrenamiento (y por tanto del ruido $\bm{\varepsilon}$). Se definen:
\[
  \text{Sesgo}[\hat{f}(\bm{x})]
  = \mathbb{E}\bigl[f(\bm{x};\hat{\bm{\theta}})\bigr] - g(\bm{x}),
\]
\[
  \text{Varianza}[\hat{f}(\bm{x})]
  = \mathbb{E}\Bigl[\bigl(f(\bm{x};\hat{\bm{\theta}}) -
    \mathbb{E}[f(\bm{x};\hat{\bm{\theta}})]\bigr)^2\Bigr].
\]
El sesgo mide cuánto se aleja la predicción media del valor verdadero. La varianza
mide cuánto fluctúa la predicción al cambiar la realización del ruido en los datos
de entrenamiento.
\end{definicion}

\begin{teorema}{Descomposición sesgo-varianza}
\label{thm:biasvar}
Sea $\bm{x}_0 \in \mathbb{R}^n$ un punto de evaluación fijo. Bajo el modelo de
ruido de la Definición~\ref{def:ruido}, el error cuadrático esperado en $\bm{x}_0$
se descompone como:
\[
  \mathbb{E}\Bigl[\bigl(f(\bm{x}_0;\hat{\bm{\theta}}) - y_0\bigr)^2\Bigr]
  \;=\;
  \underbrace{\text{Sesgo}[\hat{f}(\bm{x}_0)]^2}_{\text{error sistemático}}
  \;+\;
  \underbrace{\text{Varianza}[\hat{f}(\bm{x}_0)]}_{\text{error por ruido en train}}
  \;+\;
  \underbrace{\sigma_\varepsilon^2}_{\text{ruido irreducible}},
\]
donde $y_0 = g(\bm{x}_0) + \varepsilon_0$ con $\varepsilon_0$ independiente de
$\mathcal{D}_{\text{train}}$ y con las mismas propiedades que $\varepsilon_i$.
\end{teorema}

\begin{proof}
Denotemos $\hat{f}_0 = f(\bm{x}_0;\hat{\bm{\theta}})$ y $\bar{f}_0 =
\mathbb{E}[\hat{f}_0]$. Escribimos:
\begin{align*}
  \hat{f}_0 - y_0
  &= \hat{f}_0 - g(\bm{x}_0) - \varepsilon_0 \\
  &= \bigl(\hat{f}_0 - \bar{f}_0\bigr)
   + \bigl(\bar{f}_0 - g(\bm{x}_0)\bigr)
   - \varepsilon_0.
\end{align*}
Elevando al cuadrado y tomando esperanza, y usando que los tres términos son
no correlacionados entre sí:

\begin{itemize}
  \item $\hat{f}_0 - \bar{f}_0$ depende de $\bm{\varepsilon}$ (el ruido en
        $\mathcal{D}_{\text{train}}$) con media cero.
  \item $\bar{f}_0 - g(\bm{x}_0)$ es determinista (no depende de $\bm{\varepsilon}$).
  \item $\varepsilon_0$ es independiente de $\mathcal{D}_{\text{train}}$ con media cero.
\end{itemize}

Por tanto, los productos cruzados tienen esperanza cero, y:
\begin{align*}
  \mathbb{E}\bigl[(\hat{f}_0 - y_0)^2\bigr]
  &= \mathbb{E}\bigl[(\hat{f}_0 - \bar{f}_0)^2\bigr]
   + \bigl(\bar{f}_0 - g(\bm{x}_0)\bigr)^2
   + \mathbb{E}[\varepsilon_0^2] \\
  &= \text{Varianza}[\hat{f}(\bm{x}_0)]
   + \text{Sesgo}[\hat{f}(\bm{x}_0)]^2
   + \sigma_\varepsilon^2.
\end{align*}
\end{proof}

\begin{observacion}{Interpretación del Teorema~\ref{thm:biasvar}}
Los tres términos tienen naturaleza diferente y se comportan de forma distinta
al variar la capacidad $p$:
\begin{itemize}
  \item El \textbf{sesgo} decrece al aumentar $p$: modelos más complejos tienen más
        libertad para aproximar $g$. Para $p$ suficientemente grande, el Teorema de
        Aproximación Universal garantiza que el sesgo puede hacerse arbitrariamente
        pequeño.
  \item La \textbf{varianza} crece al aumentar $p$: modelos con más parámetros son
        más sensibles a las fluctuaciones del ruido en los datos de entrenamiento.
        Cada realización del ruido produce un $\hat{\bm{\theta}}$ diferente, y la
        dispersión de esas realizaciones define la varianza.
  \item El \textbf{ruido irreducible} $\sigma_\varepsilon^2$ no depende de $p$: es
        una propiedad del proceso de medición, no del modelo.
\end{itemize}
El mínimo del error esperado total se alcanza en el valor de $p$ donde el balance
entre sesgo y varianza es óptimo. La regularización es el mecanismo para controlar
ese balance de forma continua, sin necesidad de discretizar $p$.
\end{observacion}

\subsection{El dilema sesgo-varianza en el caso lineal}
\label{subsec:dilemma_lineal}

\dificultad{3}{2}

La descomposición del Teorema~\ref{thm:biasvar} es general, pero en el caso lineal
puede hacerse explícita. Sea la familia polinomial de grado $d$ con pérdida MSE, de
modo que $\hat{\bm{\theta}} = \bm{V}^\dagger\bm{y}$ (Proposición~\ref{prop:normales}).

Bajo el modelo $\bm{y} = \bm{V}\bm{\theta}_{\text{real}} + \bm{\varepsilon}$, el
estimador es:
\[
  \hat{\bm{\theta}} = \bm{V}^\dagger(\bm{V}\bm{\theta}_{\text{real}} + \bm{\varepsilon})
  = \bm{V}^\dagger\bm{V}\,\bm{\theta}_{\text{real}} + \bm{V}^\dagger\bm{\varepsilon}.
\]
Tomando esperanza (recordando que $\mathbb{E}[\bm{\varepsilon}] = \bm{0}$):
\[
  \mathbb{E}[\hat{\bm{\theta}}] = \bm{V}^\dagger\bm{V}\,\bm{\theta}_{\text{real}}.
\]
Cuando $\bm{V}$ tiene rango completo, $\bm{V}^\dagger\bm{V} = \bm{I}$ y el
estimador es \textbf{insesgado}: $\mathbb{E}[\hat{\bm{\theta}}] = \bm{\theta}_{\text{real}}$.
En cambio, si el modelo es mal especificado (la función real $g$ no es un polinomio
de grado $d$), el término $\bm{V}^\dagger\bm{V}\,\bm{\theta}_{\text{real}}$ introduce
sesgo sistemático, incluso con ruido cero.

La varianza del estimador es:
\[
  \mathbb{E}\bigl[\|\hat{\bm{\theta}} - \mathbb{E}[\hat{\bm{\theta}}]\|_2^2\bigr]
  = \mathbb{E}\bigl[\|\bm{V}^\dagger\bm{\varepsilon}\|_2^2\bigr]
  = \sigma_\varepsilon^2\, \|\bm{V}^\dagger\|_F^2
  = \sigma_\varepsilon^2 \sum_{i=1}^p \frac{1}{\sigma_i^2},
\]
donde $\sigma_i$ son los valores singulares de $\bm{V}$ y $\|\cdot\|_F$ es la
norma de Frobenius. Este resultado muestra que la varianza crece al aumentar $p$
(más parámetros), y se dispara cuando algún $\sigma_i \to 0$ (sistema mal condicionado).
La regularización, que introducimos a continuación, actúa precisamente sobre ese
denominador.

%───────────────────────────────────────────────────────────────────────────────
\section{Regularización de Tikhonov}
\label{sec:tikhonov}
%───────────────────────────────────────────────────────────────────────────────

\dificultad{3}{3}

La regularización modifica el problema de minimización añadiendo un término que
penaliza la complejidad del modelo. La forma más clásica y matemáticamente tractable
es la regularización de Tikhonov, conocida en ML como \textit{ridge regression} o
regularización $\ell_2$.

\begin{definicion}{Problema de minimización regularizado}
\label{def:tikhonov}
Dado un parámetro de regularización $\lambda > 0$, el \textbf{problema de
minimización regularizado con norma $\ell_2$} es:
\[
  \bm{\theta}^*_\lambda \;=\; \argmin_{\bm{\theta} \in \mathbb{R}^p}
  \Bigl\{
    \mathcal{L}(\bm{\theta}) + \lambda\|\bm{\theta}\|_2^2
  \Bigr\}
  \;=\; \argmin_{\bm{\theta} \in \mathbb{R}^p}
  \Bigl\{
    \frac{1}{N}\|\bm{V}\bm{\theta} - \bm{y}\|_2^2 + \lambda\|\bm{\theta}\|_2^2
  \Bigr\}.
\]
El escalar $\lambda > 0$ se denomina \textbf{parámetro de regularización} o
\textbf{hiperparámetro}.
\end{definicion}

\begin{observacion}{Interpretación del término de regularización}
El término $\lambda\|\bm{\theta}\|_2^2$ penaliza parámetros de magnitud grande.
Desde el punto de vista de la aproximación, parámetros grandes corresponden a
funciones que oscilan rápidamente (como el polinomio interpolante de alto grado).
La regularización impone una preferencia por soluciones más suaves. El parámetro
$\lambda$ controla la intensidad de esa preferencia: $\lambda \to 0$ recupera el
problema no regularizado, y $\lambda \to \infty$ fuerza $\bm{\theta}^*_\lambda \to \bm{0}$.
\end{observacion}

\begin{proposicion}{Solución del problema regularizado}
\label{prop:sol_tikhonov}
El minimizador del problema de la Definición~\ref{def:tikhonov} existe, es único
para todo $\lambda > 0$, y viene dado por:
\[
  \bm{\theta}^*_\lambda \;=\; \bigl(\bm{V}\T\bm{V} + N\lambda\bm{I}\bigr)^{-1}\bm{V}\T\bm{y}.
\]
\end{proposicion}

\begin{proof}
La función objetivo $F(\bm{\theta}) = \frac{1}{N}\|\bm{V}\bm{\theta}-\bm{y}\|_2^2
+ \lambda\|\bm{\theta}\|_2^2$ es suma de dos funciones convexas y diferenciables,
luego es convexa y diferenciable. Su gradiente es:
\[
  \nabla_{\bm{\theta}} F(\bm{\theta})
  = \frac{2}{N}\bm{V}\T(\bm{V}\bm{\theta} - \bm{y}) + 2\lambda\bm{\theta}
  = \frac{2}{N}\bigl[(\bm{V}\T\bm{V} + N\lambda\bm{I})\bm{\theta} - \bm{V}\T\bm{y}\bigr].
\]
Igualando a cero:
\[
  (\bm{V}\T\bm{V} + N\lambda\bm{I})\bm{\theta} = \bm{V}\T\bm{y}.
\]
La matriz $\bm{V}\T\bm{V} + N\lambda\bm{I}$ es definida positiva para todo
$\lambda > 0$: si $\bm{V}\T\bm{V}$ es semidefinida positiva (siempre) con
autovalores $\mu_i \geq 0$, entonces $\bm{V}\T\bm{V} + N\lambda\bm{I}$ tiene
autovalores $\mu_i + N\lambda > 0$. Por tanto, la matriz es invertible y la
solución es única.
\end{proof}

\begin{fisicaconexion}{Regularización en problemas inversos de física}
\label{obs:inv_problemas}
La Proposición~\ref{prop:sol_tikhonov} es la solución estándar a una clase amplia
de problemas inversos en física, donde se busca reconstruir una causa a partir de
sus efectos observados con ruido. En tomografía computarizada, se invierte la
transformada de Radon discretizada para reconstruir una imagen de densidad a partir
de proyecciones de rayos X: la matriz $\bm{V}$ es la matriz de proyección y
$\bm{y}$ son los detectores. En espectroscopía, se invierte la función de aparato
del instrumento para deconvolucionar el espectro medido y recuperar el espectro
real: el ruido en el espectrómetro produce oscilaciones espurias si no se regulariza.
En geofísica, se invierte el modelo de velocidades sísmicas a partir de tiempos de
llegada de ondas: el sistema está fuertemente subdeterminado y la regularización
impone suavidad geológica. En todos estos casos, la inversión directa
$\bm{V}^{-1}\bm{y}$ amplifica el ruido porque $\bm{V}$ es mal condicionada. La
solución $(\bm{V}\T\bm{V} + N\lambda\bm{I})^{-1}\bm{V}\T\bm{y}$ estabiliza
la inversión a costa de un sesgo controlado por $\lambda$.
\end{fisicaconexion}

\subsection{Efecto de la regularización sobre el sesgo y la varianza}
\label{subsec:sesgvar_tikhonov}

\dificultad{3}{2}

Bajo el modelo $\bm{y} = \bm{V}\bm{\theta}_{\text{real}} + \bm{\varepsilon}$, el
estimador regularizado es:
\[
  \hat{\bm{\theta}}_\lambda
  = \bm{H}_\lambda \bm{V}\T(\bm{V}\bm{\theta}_{\text{real}} + \bm{\varepsilon})
  = \bm{H}_\lambda \bm{V}\T\bm{V}\,\bm{\theta}_{\text{real}} + \bm{H}_\lambda\bm{V}\T\bm{\varepsilon},
\]
donde $\bm{H}_\lambda = (\bm{V}\T\bm{V} + N\lambda\bm{I})^{-1}$ es la matriz de
resolución.

\noindent\textbf{Sesgo del estimador regularizado:}
\[
  \text{Sesgo}[\hat{\bm{\theta}}_\lambda]
  = \mathbb{E}[\hat{\bm{\theta}}_\lambda] - \bm{\theta}_{\text{real}}
  = (\bm{H}_\lambda \bm{V}\T\bm{V} - \bm{I})\,\bm{\theta}_{\text{real}}
  = -N\lambda\,\bm{H}_\lambda\,\bm{\theta}_{\text{real}}.
\]
El sesgo es cero solo cuando $\lambda = 0$ (sin regularización). Para $\lambda > 0$,
la regularización introduce un sesgo proporcional a $\lambda$.

\noindent\textbf{Varianza del estimador regularizado:}
\[
  \text{Varianza}[\hat{\bm{\theta}}_\lambda]
  = \sigma_\varepsilon^2\,\|\bm{H}_\lambda\bm{V}\T\|_F^2
  = \sigma_\varepsilon^2 \sum_{i=1}^p \frac{\sigma_i^2}{(\sigma_i^2 + N\lambda)^2},
\]
donde $\sigma_i$ son los valores singulares de $\bm{V}$. Para $\lambda > 0$, cada
término es menor que el correspondiente sin regularización ($\sigma_\varepsilon^2/\sigma_i^2$),
y la reducción es mayor para valores singulares pequeños.

La regularización introduce sesgo y reduce varianza de forma simultánea, en
proporciones controladas por $\lambda$. La Sección~\ref{sec:svd} completa el
análisis de esta relación mediante la SVD.

%───────────────────────────────────────────────────────────────────────────────
\section{Análisis espectral: la descomposición en valores singulares}
\label{sec:svd}
%───────────────────────────────────────────────────────────────────────────────

\dificultad{3}{3}

La descomposición en valores singulares (SVD) es el instrumento de análisis lineal
más informativo para el problema de mínimos cuadrados. Permite ver exactamente
qué direcciones del espacio de parámetros están bien determinadas por los datos
y cuáles no, y entender la regularización como un filtro espectral.

\begin{teorema}{Descomposición en valores singulares (SVD)}
\label{thm:svd}
Sea $\bm{V} \in \mathbb{R}^{N \times p}$ con $N \geq p$ y $\rank(\bm{V}) = r
\leq p$. Existen matrices:
\[
  \bm{U} \in \mathbb{R}^{N \times r},\quad
  \bm{\Sigma} \in \mathbb{R}^{r \times r},\quad
  \bm{W} \in \mathbb{R}^{p \times r},
\]
con $\bm{U}\T\bm{U} = \bm{W}\T\bm{W} = \bm{I}_r$ y $\bm{\Sigma} =
\operatorname{diag}(\sigma_1, \ldots, \sigma_r)$ con $\sigma_1 \geq \sigma_2
\geq \cdots \geq \sigma_r > 0$, tales que:
\[
  \bm{V} = \bm{U}\bm{\Sigma}\bm{W}\T.
\]
Los escalares $\sigma_i$ se denominan \textbf{valores singulares} de $\bm{V}$.
Las columnas de $\bm{U}$ se denominan \textbf{vectores singulares izquierdos}
y las columnas de $\bm{W}$ \textbf{vectores singulares derechos}.
\end{teorema}

La demostración de existencia usa el teorema espectral para matrices simétricas
aplicado a $\bm{V}\T\bm{V}$; se incluye como Ejercicio Resuelto~\ref{ej:svd_construccion}
al final del capítulo.

\subsection{La pseudoinversa revisitada}
\label{subsec:pseudoinversa_svd}

Con la SVD de $\bm{V} = \bm{U}\bm{\Sigma}\bm{W}\T$, la pseudoinversa se escribe:
\[
  \bm{V}^\dagger = \bm{W}\bm{\Sigma}^{-1}\bm{U}\T,
\]
y la solución de mínimos cuadrados es $\bm{\theta}^* = \bm{W}\bm{\Sigma}^{-1}\bm{U}\T\bm{y}$.
Expandiendo en las direcciones singulares $\bm{w}_i$ (columnas de $\bm{W}$) y
$\bm{u}_i$ (columnas de $\bm{U}$):
\[
  \bm{\theta}^* = \sum_{i=1}^r \frac{\bm{u}_i\T\bm{y}}{\sigma_i}\,\bm{w}_i.
\]
Cada término $(\bm{u}_i\T\bm{y})/\sigma_i$ amplifica el componente de $\bm{y}$ en
la dirección $\bm{u}_i$ por el factor $1/\sigma_i$. Cuando $\sigma_i$ es pequeño
---dirección mal determinada por los datos--- este factor puede ser enorme, amplificando
el ruido contenido en $\bm{y}$. El número de condicionamiento $\kappa(\bm{V}) =
\sigma_1/\sigma_r$ cuantifica el peor caso de esta amplificación.

\subsection{La regularización como filtro espectral}
\label{subsec:filtro}

\dificultad{3}{3}

Con la SVD, la solución regularizada de la Proposición~\ref{prop:sol_tikhonov}
admite una descomposición espectral explícita.

\begin{proposicion}{Solución regularizada en la base singular}
\label{prop:filtro_espectral}
Sea $\bm{V} = \bm{U}\bm{\Sigma}\bm{W}\T$ la SVD de rango completo de $\bm{V}$.
El minimizador regularizado se escribe:
\[
  \bm{\theta}^*_\lambda = \sum_{i=1}^p \frac{\sigma_i}{\sigma_i^2 + N\lambda}
  \,(\bm{u}_i\T\bm{y})\,\bm{w}_i.
\]
\end{proposicion}

\begin{proof}
De la Proposición~\ref{prop:sol_tikhonov}, $\bm{\theta}^*_\lambda =
(\bm{V}\T\bm{V} + N\lambda\bm{I})^{-1}\bm{V}\T\bm{y}$. Sustituyendo
$\bm{V} = \bm{U}\bm{\Sigma}\bm{W}\T$:
\[
  \bm{V}\T\bm{V} = \bm{W}\bm{\Sigma}^2\bm{W}\T,
  \qquad
  \bm{V}\T\bm{V} + N\lambda\bm{I} = \bm{W}(\bm{\Sigma}^2 + N\lambda\bm{I})\bm{W}\T,
\]
\[
  (\bm{V}\T\bm{V} + N\lambda\bm{I})^{-1}
  = \bm{W}(\bm{\Sigma}^2 + N\lambda\bm{I})^{-1}\bm{W}\T.
\]
Por tanto:
\[
  \bm{\theta}^*_\lambda
  = \bm{W}(\bm{\Sigma}^2 + N\lambda\bm{I})^{-1}\bm{W}\T \cdot \bm{W}\bm{\Sigma}\bm{U}\T\bm{y}
  = \bm{W}\,\operatorname{diag}\!\left(\frac{\sigma_i}{\sigma_i^2+N\lambda}\right)\bm{U}\T\bm{y},
\]
que en notación de suma es la expresión del enunciado.
\end{proof}

La comparación entre la solución sin regularización y la regularizada es
inmediata:
\[
  \bm{\theta}^* = \sum_{i=1}^p \frac{1}{\sigma_i}\,(\bm{u}_i\T\bm{y})\,\bm{w}_i,
  \qquad
  \bm{\theta}^*_\lambda = \sum_{i=1}^p \underbrace{\frac{\sigma_i}{\sigma_i^2 + N\lambda}}_{f_i(\lambda)}\,(\bm{u}_i\T\bm{y})\,\bm{w}_i.
\]

Los factores $f_i(\lambda) = \sigma_i/(\sigma_i^2 + N\lambda)$ son los \textbf{factores de
filtro}. Se analizan sus propiedades:

\begin{enumerate}
  \item Para $\lambda = 0$: $f_i(0) = 1/\sigma_i$. Se recupera la pseudoinversa.
  \item Para $\sigma_i \gg \sqrt{N\lambda}$ (dirección bien determinada):
        $f_i(\lambda) \approx 1/\sigma_i$. La regularización apenas afecta a
        estas componentes.
  \item Para $\sigma_i \ll \sqrt{N\lambda}$ (dirección mal determinada):
        $f_i(\lambda) \approx \sigma_i/(N\lambda) \ll 1/\sigma_i$. La regularización
        amortigua fuertemente la contribución de estas componentes.
  \item $f_i(\lambda)$ alcanza su máximo en $\sigma_i = \sqrt{N\lambda}$, con valor
        $1/(2\sqrt{N\lambda})$.
\end{enumerate}

La regularización es, en esencia, un filtro pasa-bajos en el espacio de valores
singulares: preserva las direcciones bien determinadas ($\sigma_i$ grande) y suprime
las que el ruido dominaría ($\sigma_i$ pequeño). El umbral de corte está en
$\sigma_i \sim \sqrt{N\lambda}$.

\begin{observacion}{Analogía con la suavización espectral en física}
La estructura del filtro $f_i(\lambda) = \sigma_i/(\sigma_i^2 + N\lambda)$ es
idéntica al filtro de Wiener en procesamiento de señales: $H(\omega) =
S_x(\omega)/(S_x(\omega) + S_n(\omega))$, donde $S_x$ y $S_n$ son las densidades
espectrales de señal y ruido. En ambos casos, la supresión es proporcional a la
relación señal-ruido local en el espacio de frecuencias (o valores singulares). La
regularización de Tikhonov es un filtro de Wiener con relación señal-ruido $\sigma_i^2/N\lambda$.
\end{observacion}

%───────────────────────────────────────────────────────────────────────────────
\section{Selección del parámetro de regularización}
\label{sec:lambda}
%───────────────────────────────────────────────────────────────────────────────

\dificultad{2}{2}

El valor de $\lambda$ no puede determinarse por minimización de $\mathcal{L}_{\text{train}}$
(que siempre favorece $\lambda = 0$). Se necesita un criterio externo basado en
la capacidad del modelo de generalizar a datos nuevos.

\subsection{Validación cruzada $k$-fold}
\label{subsec:cv}

\begin{definicion}{Validación cruzada $k$-fold}
\label{def:cv}
Sea $\mathcal{D}_{\text{train}}$ el conjunto de entrenamiento con $M$ elementos.
La \textbf{validación cruzada $k$-fold} consiste en:
\begin{enumerate}
  \item Dividir $\mathcal{D}_{\text{train}}$ en $k$ subconjuntos disjuntos de
        tamaño aproximadamente igual: $\mathcal{F}_1, \ldots, \mathcal{F}_k$.
  \item Para cada $j = 1, \ldots, k$: entrenar con $\mathcal{D}_{\text{train}}
        \setminus \mathcal{F}_j$ y evaluar sobre $\mathcal{F}_j$, obteniendo
        el error de validación $E_j(\lambda)$.
  \item Estimar el error de generalización para este $\lambda$ como:
        $E_{\text{CV}}(\lambda) = \frac{1}{k}\sum_{j=1}^k E_j(\lambda)$.
\end{enumerate}
El $\lambda$ óptimo es $\lambda^* = \argmin_\lambda E_{\text{CV}}(\lambda)$.
\end{definicion}

\begin{observacion}{Casos límite: leave-one-out y train-test split}
Cuando $k = M$ (cada fold tiene un solo elemento) se obtiene la validación
\textit{leave-one-out} (LOO), que es el estimador insesgado del error de
generalización pero tiene costo computacional $O(M)$ veces el de un entrenamiento.
Cuando $k = 1$ se reduce a la partición simple entrenamiento-test de la
Definición~\ref{def:particion}. En la práctica, $k = 5$ o $k = 10$ ofrece un
buen balance entre sesgo y varianza del estimador del error.
\end{observacion}

\subsection{La curva L}
\label{subsec:curva_l}

Para problemas inversos en física, donde el conjunto de datos puede ser pequeño y
la validación cruzada costosa, existe un criterio geométrico alternativo: la curva L.

\begin{definicion}{Curva L}
\label{def:curva_l}
La \textbf{curva L} es el conjunto de puntos parametrizados por $\lambda > 0$:
\[
  \bigl(\log\mathcal{L}_{\text{train}}(\bm{\theta}^*_\lambda),\;
         \log\|\bm{\theta}^*_\lambda\|_2^2\bigr) \;\subset\; \mathbb{R}^2.
\]
\end{definicion}

Esta curva tiene la forma característica de una L: para $\lambda$ pequeño, la
norma $\|\bm{\theta}^*_\lambda\|_2$ es grande y el ajuste es bueno (brazo horizontal);
para $\lambda$ grande, el ajuste es pobre pero la norma es pequeña (brazo vertical).
La esquina de la L corresponde al valor de $\lambda$ donde el modelo alcanza el mejor
compromiso entre fidelidad a los datos y regularidad de la solución. En problemas
de física, donde el número de observaciones puede ser menor que el número de parámetros,
la curva L frecuentemente da resultados similares a la validación cruzada con un
costo computacional muy inferior.

\begin{fisicaconexion}{La curva L en problemas inversos geofísicos y espectrales}
La curva L es un criterio estándar en problemas inversos de física donde el
conjunto de datos es pequeño y la validación cruzada resulta costosa. En
\textbf{tomografía sísmica}, se usa para determinar el nivel de suavidad del
modelo de velocidades del subsuelo que mejor reproduce los tiempos de llegada
observados sin sobreajustar los errores de lectura. En \textbf{deconvolución
espectral}, aparece al intentar separar líneas espectrales superpuestas: el brazo
horizontal de la L corresponde a soluciones que ajustan bien los datos pero con
perfiles espectrales ruidosos; el brazo vertical corresponde a soluciones muy
suaves que ignoran detalles reales. La esquina identifica el nivel de regularización
que preserva las características físicas del espectro sin amplificar el ruido del
detector. En \textbf{tomografía de emisión de positrones} (PET), el mismo criterio
selecciona el parámetro de penalización que controla la resolución espacial de la
imagen reconstruida frente al ruido estadístico de los fotones detectados.
\end{fisicaconexion}

%───────────────────────────────────────────────────────────────────────────────
\section{Regularización \texorpdfstring{$\ell_1$}{ell-1} y soluciones dispersas}
\label{sec:ell1}
%───────────────────────────────────────────────────────────────────────────────

\dificultad{3}{2}

La regularización de Tikhonov no es la única posibilidad. Una alternativa importante
es la penalización con la norma $\ell_1$:

\begin{definicion}{Problema regularizado con norma $\ell_1$ (LASSO)}
\label{def:lasso}
El \textbf{problema LASSO} (\textit{Least Absolute Shrinkage and Selection Operator})
es:
\[
  \bm{\theta}^*_\lambda = \argmin_{\bm{\theta} \in \mathbb{R}^p}
  \Bigl\{
    \frac{1}{N}\|\bm{V}\bm{\theta} - \bm{y}\|_2^2 + \lambda\|\bm{\theta}\|_1
  \Bigr\},
  \qquad \|\bm{\theta}\|_1 = \sum_{j=1}^p |\theta_j|.
\]
\end{definicion}

A diferencia de la regularización $\ell_2$, la penalización $\ell_1$ tiene la
propiedad de producir soluciones \textbf{dispersas}: muchos componentes de
$\bm{\theta}^*_\lambda$ son exactamente cero. La razón es geométrica.

La solución de ambos problemas regularizados puede interpretarse como el punto
de mínima pérdida empírica en la intersección con una bola de radio $c(\lambda)$:
$\ell_2$ impone $\|\bm{\theta}\|_2 \leq c$ (bola esférica) y $\ell_1$ impone
$\|\bm{\theta}\|_1 \leq c$ (bola con forma de diamante en $\mathbb{R}^2$, o
hiperoctaedro en $\mathbb{R}^p$). Las superficies de nivel de la pérdida empírica
son elipsoides centrados en $\bm{\theta}^*$. El mínimo restringido es el primer
punto de contacto entre el elipsoide que se expande y la bola de la restricción.

La bola $\ell_2$ es lisa: el contacto ocurre en un punto genérico de la frontera.
La bola $\ell_1$ tiene vértices en los ejes coordenados, donde exactamente un
conjunto de componentes es no nulo. El contacto tiene probabilidad positiva de
ocurrir en un vértice, lo que produce soluciones con exactamente $k < p$ componentes
no nulas. El parámetro $\lambda$ controla cuántas componentes sobreviven.

\begin{observacion}{Dispersidad en redes neuronales}
En redes neuronales profundas, la regularización $\ell_1$ aplicada a los pesos
de una capa puede desactivar conexiones enteras (peso exactamente cero), produciendo
redes más compactas. Sin embargo, la función $\|\bm{\theta}\|_1$ no es diferenciable
en $\theta_i = 0$, lo que complica los algoritmos de gradiente del Capítulo 5.
Este problema se resuelve con el algoritmo de gradiente proximal, que se menciona
en el Apéndice del Capítulo 5.
\end{observacion}

%───────────────────────────────────────────────────────────────────────────────
\section{Resumen del Capítulo 2}
\label{sec:resumen2}
%───────────────────────────────────────────────────────────────────────────────

La pregunta que abrió este capítulo tenía una respuesta desconcertante: el modelo
que minimiza perfectamente la pérdida empírica no es necesariamente el mejor modelo.
El sobreajuste no es un fallo del método de optimización sino una consecuencia de
que los datos de entrenamiento contienen ruido y el modelo tiene demasiados grados
de libertad para absorberlo.

El Teorema de descomposición sesgo-varianza (\ref{thm:biasvar}) formaliza ese
diagnóstico: el error esperado en datos nuevos tiene tres fuentes, y solo el sesgo
y la varianza son controlables. El sesgo y la varianza se comportan de forma opuesta
al variar la complejidad del modelo, lo que crea un dilema que no puede resolverse
simplemente aumentando la capacidad.

La regularización de Tikhonov introduce la solución: en vez de reducir discretamente
la complejidad eligiendo un valor de $p$, se añade una penalización continua
$\lambda\|\bm{\theta}\|_2^2$ que controla el balance. El análisis espectral mediante
la SVD revela la elegancia del mecanismo: la regularización es un filtro pasa-bajos
que preserva las direcciones bien determinadas por los datos y amortigua las que
serían dominadas por el ruido. La analogía con el filtro de Wiener en física de
señales es exacta.

El capítulo cierra con una observación que anticipa lo que sigue. Todo lo desarrollado
aquí es exacto porque la familia paramétrica era lineal en $\bm{\theta}$: la pérdida
empírica era una forma cuadrática, la SVD diagonalizaba el problema, y la solución
tenía forma cerrada. Cuando la familia es no lineal ---como ocurrirá a partir del
Capítulo 4--- el paisaje de $\mathcal{L}(\bm{\theta})$ deja de ser una parábola.
La geometría de ese paisaje, y la forma en que determina qué algoritmos de optimización
son viables, es el objeto del Capítulo 3.

%───────────────────────────────────────────────────────────────────────────────
\section*{Ejercicios resueltos}
\addcontentsline{toc}{section}{Ejercicios resueltos}
\label{sec:ej_resueltos2}
%───────────────────────────────────────────────────────────────────────────────

\begin{ejercicio}
\label{ej:svd_construccion}
\textbf{(Construcción de la SVD a partir del teorema espectral).}

Sea $\bm{V} \in \mathbb{R}^{N\times p}$ con $N \geq p$. Mostrar que existe una
descomposición $\bm{V} = \bm{U}\bm{\Sigma}\bm{W}\T$ con las propiedades del
Teorema~\ref{thm:svd}.

\medskip\noindent\textit{Solución.}
La matriz $\bm{V}\T\bm{V} \in \mathbb{R}^{p\times p}$ es simétrica (pues
$(\bm{V}\T\bm{V})\T = \bm{V}\T\bm{V}$) y semidefinida positiva:
$\bm{z}\T(\bm{V}\T\bm{V})\bm{z} = \|\bm{V}\bm{z}\|_2^2 \geq 0$ para todo
$\bm{z} \in \mathbb{R}^p$.

Por el teorema espectral para matrices simétricas reales, existe una
descomposición:
\[
  \bm{V}\T\bm{V} = \bm{W}\bm{\Lambda}\bm{W}\T,
\]
donde $\bm{W} \in \mathbb{R}^{p\times p}$ es ortogonal ($\bm{W}\T\bm{W} = \bm{I}$)
y $\bm{\Lambda} = \operatorname{diag}(\mu_1, \ldots, \mu_p)$ con $\mu_i \geq 0$.

Definimos $\sigma_i = \sqrt{\mu_i}$ y $\bm{\Sigma} = \operatorname{diag}(\sigma_1,\ldots,\sigma_r)$
donde $r = \rank(\bm{V})$ es el número de autovalores no nulos. Para
$i = 1, \ldots, r$ definimos:
\[
  \bm{u}_i = \frac{\bm{V}\bm{w}_i}{\sigma_i} \in \mathbb{R}^N,
\]
donde $\bm{w}_i$ es la $i$-ésima columna de $\bm{W}$. Verificamos que
$\{\bm{u}_i\}_{i=1}^r$ es un conjunto ortonormal:
\[
  \bm{u}_i\T\bm{u}_j
  = \frac{(\bm{V}\bm{w}_i)\T(\bm{V}\bm{w}_j)}{\sigma_i\sigma_j}
  = \frac{\bm{w}_i\T(\bm{V}\T\bm{V})\bm{w}_j}{\sigma_i\sigma_j}
  = \frac{\bm{w}_i\T(\mu_j\bm{w}_j)}{\sigma_i\sigma_j}
  = \frac{\mu_j\delta_{ij}}{\sigma_i\sigma_j} = \delta_{ij}.
\]
Definiendo $\bm{U}$ como la matriz de columnas $\bm{u}_1,\ldots,\bm{u}_r$,
verificamos:
\[
  \bm{U}\bm{\Sigma}\bm{W}\T
  \;=\; \sum_{i=1}^r \sigma_i\,\bm{u}_i\bm{w}_i\T
  \;=\; \sum_{i=1}^r \sigma_i\,\frac{\bm{V}\bm{w}_i}{\sigma_i}\bm{w}_i\T
  \;=\; \bm{V}\sum_{i=1}^r \bm{w}_i\bm{w}_i\T = \bm{V},
\]
donde la última igualdad usa que los vectores $\bm{w}_i$ para $i > r$
están en el núcleo de $\bm{V}$ y no contribuyen a la suma. $\square$
\end{ejercicio}

\begin{ejercicio}
\label{ej:tikhonov_numerico}
\textbf{(Efecto de la regularización: experimento numérico).}

Sea $g(x) = \sin(2\pi x)$ muestreada en $N = 15$ nodos equiespaciados en $[0,1]$
con ruido gaussiano $\varepsilon_i \sim \mathcal{N}(0, 0.1^2)$. Se ajusta una
familia polinomial de grado $d = 14$ (interpolación exacta) y se comparan tres
valores: $\lambda = 0$ (sin regularización), $\lambda = 10^{-3}$, $\lambda = 10^{-1}$.

Calcular analíticamente (usando la SVD) el número de condicionamiento de
$\bm{V}\T\bm{V}$ para nodos equiespaciados con $d = 14$, y estimar cuánto
amplifica el ruido la solución no regularizada.

\medskip\noindent\textit{Solución parcial.}
Los nodos equiespaciados $x_i = (i-1)/14$ producen una matriz de Vandermonde
$\bm{V} \in \mathbb{R}^{15 \times 15}$ cuyo número de condicionamiento crece
exponencialmente con $d$. Para $d = 14$, $\kappa(\bm{V}) \approx 10^{13}$
(verificable numéricamente). Esto significa que una perturbación de amplitud
$0.1$ en $\bm{y}$ puede producir una perturbación de amplitud
$\approx 10^{12}$ en $\bm{\theta}^*$, lo que explica las oscilaciones violentas
del polinomio interpolante.

Con $\lambda = 10^{-1}$, el factor de filtro para la dirección más inestable
($\sigma_r \approx 10^{-6}$) pasa de $1/\sigma_r \approx 10^6$ a
$\sigma_r/(\sigma_r^2 + N\lambda) \approx \sigma_r/(N\lambda) \approx 7\times10^{-7}$:
una reducción de más de 12 órdenes de magnitud en la amplificación del ruido.
La solución regularizada recupera la forma cualitativa de $\sin(2\pi x)$ a costa
de no interpolar exactamente los datos.
\end{ejercicio}

\begin{ejercicio}
\label{ej:descomp_svd}
\textbf{(SVD de una matriz pequeña).}

Sea $\bm{V} = \begin{pmatrix} 1 & 1 \\ 1 & -1 \\ 0 & 1 \end{pmatrix}$.
Calcular explícitamente la SVD $\bm{V} = \bm{U}\bm{\Sigma}\bm{W}\T$.

\medskip\noindent\textit{Solución.}
\[
  \bm{V}\T\bm{V} =
  \begin{pmatrix}1&1&0\\1&-1&1\end{pmatrix}
  \begin{pmatrix}1&1\\1&-1\\0&1\end{pmatrix}
  = \begin{pmatrix}2&0\\0&3\end{pmatrix}.
\]
Los autovalores son $\mu_1 = 3$, $\mu_2 = 2$, con autovectores
$\bm{w}_1 = (0,1)\T$ y $\bm{w}_2 = (1,0)\T$. Entonces:
$\sigma_1 = \sqrt{3}$, $\sigma_2 = \sqrt{2}$, $\bm{W} = \begin{pmatrix}0&1\\1&0\end{pmatrix}$.

Los vectores singulares izquierdos son:
\[
  \bm{u}_1 = \frac{\bm{V}\bm{w}_1}{\sigma_1}
  = \frac{1}{\sqrt{3}}\begin{pmatrix}1\\-1\\1\end{pmatrix},
  \qquad
  \bm{u}_2 = \frac{\bm{V}\bm{w}_2}{\sigma_2}
  = \frac{1}{\sqrt{2}}\begin{pmatrix}1\\1\\0\end{pmatrix}.
\]

Por tanto:
\[
  \bm{U} = \begin{pmatrix}
    1/\sqrt{3} & 1/\sqrt{2} \\
    -1/\sqrt{3} & 1/\sqrt{2} \\
    1/\sqrt{3} & 0
  \end{pmatrix},\quad
  \bm{\Sigma} = \begin{pmatrix}\sqrt{3}&0\\0&\sqrt{2}\end{pmatrix},\quad
  \bm{W} = \begin{pmatrix}0&1\\1&0\end{pmatrix}.
\]
Verificación: $\bm{U}\T\bm{U} = \bm{I}_2$, $\bm{W}\T\bm{W} = \bm{I}_2$,
$\bm{U}\bm{\Sigma}\bm{W}\T = \bm{V}$. $\square$
\end{ejercicio}

%───────────────────────────────────────────────────────────────────────────────
\section*{Ejercicios propuestos}
\addcontentsline{toc}{section}{Ejercicios propuestos}
\label{sec:ej_propuestos2}
%───────────────────────────────────────────────────────────────────────────────

\begin{ejercicio}
\textbf{(Fenómeno de Runge y partición entrenamiento-test).}
Sea $g(x) = 1/(1+25x^2)$ en $[-1,1]$.
\begin{enumerate}[label=(\alph*)]
  \item Para $N = 5, 9, 13$ nodos equiespaciados, ajuste el polinomio interpolante
        de grado $N-1$ y calcule el error máximo $\max_{x\in[-1,1]}|f^*(x)-g(x)|$
        en 200 puntos uniformes.
  \item Repita con nodos de Chebyshev $x_k = \cos\bigl(\frac{(2k-1)\pi}{2N}\bigr)$.
  \item ¿Para qué tipo de nodos el error disminuye al aumentar $N$? ¿Por qué?
        Relacione con el análisis de la sección~\ref{sec:svd}.
\end{enumerate}
\end{ejercicio}

\begin{ejercicio}
\textbf{(Descomposición sesgo-varianza explícita).}
Sea la familia lineal $f(x;\theta) = \theta x$ (una sola recta por el origen)
y los datos $y_i = 2x_i + \varepsilon_i$ con $\varepsilon_i \sim \mathcal{N}(0,\sigma^2)$
y $x_i$ fijos.
\begin{enumerate}[label=(\alph*)]
  \item Calcule $\hat{\theta} = \argmin_\theta \sum_i (y_i - \theta x_i)^2$.
        Expréselo en función de $x_i$ e $y_i$.
  \item Calcule $\mathbb{E}[\hat\theta]$ y $\text{Var}[\hat\theta]$.
  \item Calcule el sesgo y la varianza de la predicción $\hat{f}(x_0) = \hat\theta x_0$
        para un punto fijo $x_0$.
  \item Verifique el Teorema~\ref{thm:biasvar} calculando los tres términos
        explícitamente.
\end{enumerate}
\end{ejercicio}

\begin{ejercicio}
\textbf{(Curva de regularización).}
Sea $\bm{V} \in \mathbb{R}^{20 \times 5}$ una matriz de Vandermonde con nodos
equiespaciados en $[0,1]$ y $\bm{y} = \bm{V}\bm{\theta}_{\text{real}} + \bm{\varepsilon}$
con $\bm{\theta}_{\text{real}} = (1,0,-1,0,1)\T$ y $\varepsilon_i \sim \mathcal{N}(0, 0.5^2)$.
\begin{enumerate}[label=(\alph*)]
  \item Calcule la SVD de $\bm{V}$ y grafique los valores singulares.
  \item Para $\lambda \in \{10^{-4}, 10^{-3}, 10^{-2}, 10^{-1}, 1, 10\}$, calcule
        $\bm{\theta}^*_\lambda$, $E_{\text{train}}(\lambda)$, y
        $\|\bm{\theta}^*_\lambda\|_2$.
  \item Trace la curva L (log $E_{\text{train}}$ vs. log $\|\bm{\theta}^*_\lambda\|_2$)
        e identifique visualmente la esquina.
  \item Compare el $\lambda^*$ obtenido de la curva L con el obtenido por validación
        cruzada 5-fold.
\end{enumerate}
\end{ejercicio}

\begin{ejercicio}
\textbf{(LASSO vs. Ridge: interpretación geométrica).}
En $\mathbb{R}^2$ con $\bm{\theta} = (\theta_1, \theta_2)\T$, considérese una pérdida
empírica cuadrática $\mathcal{L}(\bm{\theta}) = (\theta_1 - 3)^2 + (\theta_2 - 2)^2$
(es decir, el mínimo sin restricciones está en $\bm{\theta}^* = (3,2)\T$).
\begin{enumerate}[label=(\alph*)]
  \item Dibuje las curvas de nivel de $\mathcal{L}$ y las bolas $\|\bm{\theta}\|_2
        \leq c$ y $\|\bm{\theta}\|_1 \leq c$ para $c = 1, 2, 3$.
  \item Para la restricción $\|\bm{\theta}\|_2^2 \leq c^2$, calcule analíticamente
        el minimizador restringido como función de $c$.
  \item Para la restricción $\|\bm{\theta}\|_1 \leq c$, determine si el mínimo
        restringido cae sobre un eje coordenado (solución dispersa) o en un
        punto interior de un lado del diamante.
  \item ¿Para qué valor de $c$ la solución LASSO tiene $\theta_1 = 0$?
\end{enumerate}
\end{ejercicio}

\begin{ejercicio}
\textbf{(Invarianza de la regularización $\ell_2$ bajo rotación).}
Sea $\bm{\theta}' = \bm{R}\bm{\theta}$ con $\bm{R}$ ortogonal. Mostrar que:
\begin{enumerate}[label=(\alph*)]
  \item $\|\bm{\theta}'\|_2 = \|\bm{\theta}\|_2$ (la regularización $\ell_2$ es
        invariante bajo rotación del espacio de parámetros).
  \item $\|\bm{\theta}'\|_1 \neq \|\bm{\theta}\|_1$ en general (la regularización
        $\ell_1$ no es invariante bajo rotación).
  \item Discuta la implicación: ¿la interpretación de la dispersidad en LASSO
        depende de la elección de base en el espacio de parámetros?
\end{enumerate}
\end{ejercicio}

%───────────────────────────────────────────────────────────────────────────────
\label{sec:estrella2}
\begin{estrella}{Reconstrucción del flujo neutrónico con datos ruidosos}

\textbf{El sistema físico.}
El reactor del Capítulo~\ref{ch:cap1} ahora opera en condiciones reales: los
$N$ detectores reportan $\phi_i = \phi(\bm{x}_i) + \varepsilon_i$ donde
$\varepsilon_i$ es ruido electrónico con varianza $\sigma_\varepsilon^2$. El
número de funciones base RBF a usar no está fijado: el operador puede elegir
$p$ entre $1$ y $N$.

\textbf{El dilema operacional.}
Con $p = N$ (interpolación exacta, $\lambda = 0$): la distribución de potencia
reconstruida pasa exactamente por cada lectura, pero oscila entre detectores
---el fenómeno de Runge en versión nuclear. Con $p = 1$ (un solo modo, $\lambda \to \infty$):
la reconstrucción es suave pero ignora la heterogeneidad real del flujo. El
parámetro de regularización $\lambda$ controla este balance de forma continua.

\textbf{El criterio de selección.}
En el reactor, el nivel de ruido $\sigma_\varepsilon^2$ es conocido (calibración
del sistema de instrumentación). La selección óptima de $\lambda$ puede hacerse
igualando la norma del término de sesgo con la norma del término de varianza:
\[
  N^2\lambda^2 \|\bm{H}_\lambda\,\bm{\theta}_{\text{real}}\|_2^2
  = \sigma_\varepsilon^2 \sum_{i=1}^p \frac{\sigma_i^2}{(\sigma_i^2 + N\lambda)^2}.
\]
Cuando $\sigma_\varepsilon^2$ no se conoce con precisión, la curva L
(Definición~\ref{def:curva_l}) proporciona el $\lambda^*$ con costo computacional
$O(p^3)$ más la SVD, frente a $O(kp^3)$ de la validación cruzada $k$-fold.

\textbf{Conexión con el resto del libro.}
La SVD de la matriz de diseño RBF diagonaliza el problema de reconstrucción:
los vectores singulares derechos $\bm{w}_i$ son los modos espaciales del flujo
que los detectores pueden ver, y los valores singulares $\sigma_i$ cuantifican
con qué precisión. Los modos con $\sigma_i \ll \sqrt{N\lambda}$ son suprimidos
por la regularización ---exactamente los modos de alta frecuencia espacial que
el ruido de los detectores no puede distinguir de la señal. En el Capítulo~9,
este análisis modal se reformula en términos de las funciones propias de la
ecuación de difusión, que tienen simetría discreta $O_h$ heredada de la geometría
cúbica del reactor.
\end{estrella}

%───────────────────────────────────────────────────────────────────────────────
\section*{Ejercicios computacionales}
\addcontentsline{toc}{section}{Ejercicios computacionales}
%───────────────────────────────────────────────────────────────────────────────

\begin{ejerciciocomp}{Regularización de Tikhonov y curva L para flujo neutrónico}
\textbf{Objetivo:} implementar la regularización de Tikhonov vía SVD y seleccionar
$\lambda$ por curva L para el problema de reconstrucción del flujo del reactor.

\medskip\noindent\textbf{Datos:}
$N = 12$ detectores en posiciones $x_i = i \cdot L/(N+1)$ con $L = 100$ cm.
Lecturas sintéticas: $\phi_i = \sin(\pi x_i/L) + \varepsilon_i$ con
$\varepsilon_i \sim \mathcal{N}(0, 0.02^2)$ (ruido del 2\% del máximo del flujo).
Familia: $p = 12$ funciones RBF gaussianas con centros $c_j = x_j$ y $\sigma = 12$ cm.

\medskip\noindent\textbf{Algoritmo:}
\begin{enumerate}
  \item Construir la matriz de diseño $\bm{V} \in \mathbb{R}^{12\times 12}$ con
        $V_{ij} = \exp\bigl(-\|x_i - c_j\|^2/(2\sigma^2)\bigr)$.
  \item Calcular la SVD $\bm{V} = \bm{U}\bm{\Sigma}\bm{W}\T$ y reportar el
        número de condicionamiento $\kappa(\bm{V})$.
  \item Para $\lambda \in \{10^{-8}, 10^{-7}, \ldots, 10^{0}\}$ (9 valores),
        calcular $\bm{\theta}^*_\lambda$ usando la fórmula espectral de la
        Proposición~\ref{prop:filtro_espectral}. Calcular
        $E_{\text{train}}(\lambda)$ y $\|\bm{\theta}^*_\lambda\|_2$.
  \item Trazar la curva L en escala log-log e identificar visualmente
        la esquina. Reportar el $\lambda^*$ correspondiente.
  \item Graficar $\phi_\lambda^*(x)$ en 200 puntos uniformes en $[0, L]$
        para $\lambda = 0$ (pseudoinversa), $\lambda = \lambda^*$ y
        $\lambda = 10^{-1}$ (sobresuavizado).
\end{enumerate}

\medskip\noindent\textbf{Verificación:}
Para $\lambda = \lambda^*$ (típicamente $\approx 10^{-5}$ a $10^{-4}$):
el error RMS entre $\phi_\lambda^*(x)$ y $\sin(\pi x/L)$ en 200 puntos debe
ser $< 0.05$ (unidades del flujo máximo). Para $\lambda = 0$, el error RMS
debe ser $> 0.5$ debido al sobreajuste.

\medskip\noindent\textbf{Extensión:}
Repetir con $N = 6$ detectores (sistema subdeterminado, $p > N$). En este caso,
la pseudoinversa da infinitas soluciones; la regularización selecciona la de
norma mínima. ¿Cómo cambia la curva L cuando $p > N$?
\end{ejerciciocomp}

\begin{ejerciciocomp}{Validación cruzada $k$-fold para selección de $\lambda$}
\textbf{Objetivo:} comparar la curva L con la validación cruzada 5-fold como
criterios de selección de $\lambda$ en el mismo problema de flujo.

\medskip\noindent\textbf{Datos:}
Misma configuración que el ejercicio anterior con $N = 20$ detectores y
$\sigma_\varepsilon = 0.03$.

\medskip\noindent\textbf{Algoritmo:}
\begin{enumerate}
  \item Implementar la validación cruzada 5-fold: dividir los 20 detectores
        en 5 grupos de 4. Para cada fold $j$: entrenar con los 16 detectores
        restantes, evaluar el MSE sobre los 4 del fold. Promediar los 5 errores.
  \item Repetir para los mismos 9 valores de $\lambda$ del ejercicio anterior.
  \item Comparar el $\lambda^*_{\text{CV}}$ con el $\lambda^*_{\text{L}}$
        obtenido por curva L.
  \item Calcular el tiempo de cómputo de cada método.
\end{enumerate}

\medskip\noindent\textbf{Verificación:}
El cociente $\lambda^*_{\text{CV}} / \lambda^*_{\text{L}}$ debe estar en el
intervalo $[0.1, 10]$: ambos métodos deben dar órdenes de magnitud similares.
El tiempo de la validación cruzada debe ser aproximadamente $5\times$ el de
la curva L.

\medskip\noindent\textbf{Extensión:}
Repetir el experimento 100 veces con distintas realizaciones del ruido y
calcular la distribución de $\lambda^*_{\text{CV}}$ y $\lambda^*_{\text{L}}$.
¿Cuál tiene menor varianza? ¿Cuál tiene menor sesgo respecto al $\lambda$
que minimiza el error verdadero (conocido por la función sintética)?
\end{ejerciciocomp}

%───────────────────────────────────────────────────────────────────────────────
\begin{notasbib}
La regularización de Tikhonov fue formulada por Andrei Nikolaevich Tikhonov en
1963 en el contexto de ecuaciones integrales de Fredholm de primer tipo, que son
el prototipo de problema inverso mal condicionado~\cite{tikhonov1963solution}.
En el contexto de mínimos cuadrados estadísticos, la misma solución fue propuesta
independientemente por Hoerl y Kennard (1970) bajo el nombre de \textit{ridge
regression}. La conexión entre ambas formulaciones —probabilística y variacional—
se hace explícita en la interpretación bayesiana del Capítulo~7.

La descomposición sesgo-varianza en la forma del Teorema~\ref{thm:biasvar} es
estándar; una presentación rigurosa con condiciones mínimas sobre el ruido se
encuentra en Hastie, Tibshirani y Friedman~\cite{hastie2009elements}. La extensión
al caso no lineal (redes neuronales) es más sutil porque el estimador depende no
linealmente de los datos; el análisis correspondiente usa herramientas de funciones
de influencia y bootstrap.

La SVD como herramienta para entender la regularización es el tema central del
libro clásico de Hansen~\cite{hansen1998rank}, donde se desarrollan también el
criterio de la curva L y el criterio de la validación cruzada generalizada (GCV).
La analogía entre la regularización de Tikhonov y el filtro de Wiener aparece
explícitamente en el Capítulo~5 de ese texto.

El despliegue espectral neutrónico (\textit{neutron spectrum unfolding}) como
aplicación de la regularización de Tikhonov en física nuclear está documentado en
McElroy et al.\ (SAND-II, 1967) y en Reginatto~\cite{reginatto2010overview}.
La curva L como criterio de selección de $\lambda$ fue analizada sistemáticamente
por Hansen~\cite{hansen1998rank}, donde se presentan también aplicaciones a
tomografía sísmica, deconvolución espectral y tomografía de emisión de positrones.

El LASSO fue introducido por Tibshirani~\cite{tibshirani1996regression} en 1996.
La interpretación geométrica mediante la bola $\ell_1$ es la presentación estándar;
la prueba de que la solución es dispersa con probabilidad uno en el caso continuo
(cuando el mínimo no regularizado cae fuera de la bola) se demuestra en el
Capítulo~3 de Wainwright~\cite{wainwright2019high}. El algoritmo de gradiente
proximal para resolver el LASSO se trata en el Apéndice del Capítulo~5.
\end{notasbib}